% =========================================================
\section{Semaine 9 : EDOs scalaires et Systèmes}
% =========================================================

\subsection{Méthodes à 1 pas et Erreur}

\textbf{Rappel :} On considère le problème continu
\begin{equation}
\begin{cases}
    \frac{du}{dt}(t) = f(t, u(t)), \quad t \in (0, T] \\
    u(0) = u_0
\end{cases} \label{eq:edo}
\end{equation}

On partitionne l'intervalle de temps $[0, T]$. Soit $\Delta t = \frac{T}{N}$. On définit les instants discrets :
$$0 = t_0 < t_1 < \dots < t_n = t_0 + n\Delta t < \dots < t_N = T$$

\begin{definition}[Méthode à 1 pas]
Une méthode à 1 pas pour approximer la solution de \eqref{eq:edo} est de la forme :
\begin{equation}
\frac{u_{n+1} - u_n}{\Delta t} = \Phi_f(u_n, u_{n+1}, t_n, \Delta t)
\end{equation}
\end{definition}

\begin{explication}
Dans le cas d'une méthode explicite, la fonction d'incrément $\Phi_f$ ne dépend pas de $u_{n+1}$.
\end{explication}

\begin{algorithme}[Fonctions d'incrément classiques]
    \centering
    \renewcommand{\arraystretch}{1.2}
    \begin{tabular}{ll}
        \toprule
        \textbf{Méthode} & $\mathbf{\Phi_f}$ \\
        \midrule
        EP (Euler Progressif) & $f(t_n, u_n)$ \\
        ER (Euler Rétrograde) & $f(t_{n+1}, u_{n+1})$ \\
        CN (Crank-Nicolson) & $\frac{1}{2}\big(f(t_n, u_n) + f(t_{n+1}, u_{n+1})\big)$ \\
        Heun & $\frac{1}{2}\big(f(t_n, u_n) + f(t_{n+1}, u_n + \Delta t f(t_n, u_n))\big)$ \\
        RK4 & $\frac{1}{6}(k_1 + 2k_2 + 2k_3 + k_4)$ \\
        \bottomrule
    \end{tabular}
\end{algorithme}

\begin{definition}[Erreur de troncature locale]
L'erreur de troncature locale au temps $t_{n+1}$ est définie par :
\begin{equation}
\tau_{n+1} = \frac{u(t_{n+1}) - u(t_n)}{\Delta t} - \Phi_f(u(t_n), u(t_{n+1}), t_n, \Delta t)
\end{equation}
\end{definition}

\begin{theoreme}[Consistance et Convergence]
Si $f$ est lipschitzienne par rapport à son second argument et si la méthode est consistante à l'ordre $p$ ($\tau(T) = \mathcal{O}(\Delta t^p)$), alors elle est convergente à l'ordre $p$ ($e(T) = \mathcal{O}(\Delta t^p)$).
\end{theoreme}

\begin{exemple}[Euler Progressif]
Pour EP, par développement de Taylor :
$$ \tau_{n+1} \le \frac{\Delta t}{2} \max_{t \in [0, T]} |u''(t)| $$
La méthode est donc consistante et convergente à l'ordre 1.
\end{exemple}

\subsection{Stabilité}

\begin{explication}[Choix des méthodes implicites]
Considérons $\frac{du}{dt} = -\lambda u$ avec $\lambda > 0$.
\begin{itemize}
    \item \textbf{EP} : $u_{n+1} = (1 - \lambda \Delta t)u_n$. Stable ssi $\Delta t < \frac{2}{\lambda}$ (Conditionnelle).
    \item \textbf{ER} : $u_{n+1} = \frac{1}{1 + \lambda \Delta t} u_n$. Stable $\forall \Delta t > 0$ (Inconditionnelle).
\end{itemize}
\end{explication}

\subsection{Systèmes d'EDOs}

Pour $\vec{u}(t) \in \mathbb{R}^k$, on utilise les mêmes schémas sous forme vectorielle. La stabilité pour $\frac{d\vec{u}}{dt} = -A\vec{u}$ dépend des valeurs propres $\lambda_i$ de $A$. EP est stable si $|1 - \Delta t \lambda_i| < 1$ pour toutes les valeurs propres.

\subsection{EDOs d'ordre supérieur}

\begin{algorithme}[Réduction à l'ordre 1]
On transforme $\frac{d^2u}{dt^2} = f(t, u, u')$ en un système d'ordre 1 en posant $v = u'$ :
$$ \frac{d}{dt} \begin{pmatrix} u \\ v \end{pmatrix} = \begin{pmatrix} v \\ f(t, u, v) \end{pmatrix} $$
\end{algorithme}
